\chapter{2016年浙江卷（文）}

\section{单选题}

\begin{problem}
已知全集 \(U=\{1,2,3,4,5,6\}\)，集合 \(P=\{1,3,5\}\)，\(Q=\{1,2,4\}\)，则 \((\complement_U P)\cup Q=\)（\quad）

\choices
    {\(\{1\}\)}
    {\(\{3,5\}\)}
    {\(\{1,2,4,6\}\)}
    {\(\{1,2,3,4,5\}\)}
\end{problem}

\begin{answer}
C
\end{answer}

\begin{problem}
已知互相垂直的平面 \(\alpha\)，\(\beta\) 交于直线 \(l\)，若直线 \(m\)，\(n\) 满足 \(m\parallel\alpha\)，\(n\perp\beta\)，则（\quad）

\choices
    {\(m\parallel l\)}
    {\(m\parallel n\)}
    {\(n\perp l\)}
    {\(m\perp n\)}
\end{problem}

\begin{answer}
C
\end{answer}

\begin{problem}
函数 \(y=\sin x^2\) 的图象是（\quad）

\choices
    {\choicebitmap[width=0.15\paperwidth]{img/2016/zhejiang_liberal/q3_optA.png}}
    {\choicebitmap[width=0.15\paperwidth]{img/2016/zhejiang_liberal/q3_optB.png}}
    {\choicebitmap[width=0.15\paperwidth]{img/2016/zhejiang_liberal/q3_optC.png}}
    {\choicebitmap[width=0.15\paperwidth]{img/2016/zhejiang_liberal/q3_optD.png}}
\end{problem}

\begin{answer}
D
\end{answer}

\begin{problem}
若平面区域
\(\begin{cases}
x+y-3\ge 0,\\
2x-y-3\le 0,\\
x-2y+3\ge 0
\end{cases}\)
夹在两条斜率为 \(1\) 的平行直线之间，则这两条平行直线间的距离的最小值是（\quad）

\choices
    {\(\dfrac{3\sqrt5}{5}\)}
    {\(\sqrt2\)}
    {\(\dfrac{3\sqrt2}{2}\)}
    {\(\sqrt5\)}
\end{problem}

\begin{answer}
B
\end{answer}

\begin{problem}
已知 \(a,b>0\)，且 \(a\ne1\)，\(b\ne1\)，若 \(\log_a b>1\)，则（\quad）

\choices
    {\((a-1)(b-1)<0\)}
    {\((a-1)(a-b)>0\)}
    {\((b-1)(b-a)<0\)}
    {\((b-1)(b-a)>0\)}
\end{problem}

\begin{answer}
D
\end{answer}

\begin{problem}
已知函数 \(f(x)=x^2+bx\)，则“\(b<0\)”是“\(f(f(x))\) 的最小值与 \(f(x)\) 的最小值相等”的（\quad）

\choices
    {\text{充分不必要条件}}
    {\text{必要不充分条件}}
    {\text{充分必要条件}}
    {\text{既不充分也不必要条件}}
\end{problem}

\begin{answer}
A
\end{answer}

\begin{problem}
已知函数 \(f(x)\) 满足：\(f(x)\ge |x|\) 且 \(f(x)\ge 2^x\)，\(x\in\R\). 则（\quad）

\choices
    {若 \(f(a)\le |b|\)，则 \(a\le b\)}
    {若 \(f(a)\le 2^b\)，则 \(a\le b\)}
    {若 \(f(a)\ge |b|\)，则 \(a\ge b\)}
    {若 \(f(a)\ge 2^b\)，则 \(a\ge b\)}
\end{problem}

\begin{answer}
B
\end{answer}

\begin{problem}
如图，点列 \(\{A_n\}\)，\(\{B_n\}\) 分别在某锐角的两边上，且 \(|A_nA_{n+1}|=|A_{n+1}A_{n+2}|\)，\(A_n\ne A_{n+2}\)，\(n\in\N^*\)，\(|B_nB_{n+1}|=|B_{n+1}B_{n+2}|\)，\(B_n\ne B_{n+2}\)，\(n\in\N^*\)（\(P\ne Q\) 表示点 \(P\) 与点 \(Q\) 不重合）. 若 \(d_n=|A_nB_n|\)，\(S_n\) 为 \(\triangle A_nB_nB_{n+1}\) 的面积，则（\quad）

\bitmapinlinefigure[5.2cm][max width=5.2cm]{img/2016/zhejiang_liberal/q8_fig1.png}

\choices
    {\(\{S_n\}\) 是等差数列}
    {\(\{S_n^2\}\) 是等差数列}
    {\(\{d_n\}\) 是等差数列}
    {\(\{d_n^2\}\) 是等差数列}
\end{problem}

\begin{answer}
A
\end{answer}

\section{填空题}

\begin{problem}
某几何体的三视图如图所示（单位：\(\mathrm{cm}\)），则该几何体的表面积是\fillinblank{}\(\mathrm{cm}^2\)，体积是\fillinblank{}\(\mathrm{cm}^3\).

\bitmapfigure[float=inline,max width=5.0cm]{img/2016/zhejiang_liberal/q9_fig1.png}
\end{problem}

\begin{answer}
\(80\)，\(40\)
\end{answer}

\begin{problem}
已知 \(a\in\R\)，方程 \(a^2x^2+(a+2)y^2+4x+8y+5a=0\) 表示圆，则圆心坐标是\fillinblank{}，半径是\fillinblank{}.
\end{problem}

\begin{answer}
\((-2,-4)\)，\(5\)
\end{answer}

\begin{problem}
已知 \(2\cos^2x+\sin2x=A\sin(\omega x+\varphi)+B\)（\(A>0\)），则 \(A=\fillinblank{}\)，\(B=\fillinblank{}\).
\end{problem}

\begin{answer}
\(\sqrt2\)，\(1\)
\end{answer}

\begin{problem}
设函数 \(f(x)=x^3+3x^2+1\)，已知 \(a\ne0\)，且 \(f(x)-f(a)=(x-b)(x-a)^2\)，\(x\in\R\)，则实数 \(a=\fillinblank{}\)，\(b=\fillinblank{}\).
\end{problem}

\begin{answer}
\(-2\)，\(1\)
\end{answer}

\begin{problem}
设双曲线 \(x^2-\dfrac{y^2}{3}=1\) 的左，右焦点分别为 \(F_1\)，\(F_2\)，若点 \(P\) 在双曲线上，且 \(\triangle F_1PF_2\) 为锐角三角形，则 \(|PF_1|+|PF_2|\) 的取值范围是\fillinblank{}.
\end{problem}

\begin{answer}
\(\left(2\sqrt7,8\right)\)
\end{answer}

\begin{problem}
如图，已知平面四边形 \(ABCD\)，\(AB=BC=3\)，\(CD=1\)，\(AD=\sqrt5\)，\(\angle ADC=90^\circ\)，沿直线 \(AC\) 将 \(\triangle ACD\) 翻折成 \(\triangle ACD'\)，直线 \(AC\) 与 \(BD'\) 所成角的余弦的最大值是\fillinblank{}.

\bitmapfigure[float=inline,max width=4.2cm]{img/2016/zhejiang_liberal/q14_fig1.png}
\end{problem}

\begin{answer}
\(\dfrac{\sqrt6}{6}\)
\end{answer}

\begin{problem}
已知平面向量 \(\bs a\)，\(\bs b\)，\(|\bs a|=1\)，\(|\bs b|=2\)，\(\bs a\cdot\bs b=1\)，若 \(\bs e\) 为平面单位向量，则 \(|\bs a\cdot\bs e|+|\bs b\cdot\bs e|\) 的最大值是\fillinblank{}.
\end{problem}

\begin{answer}
\(\sqrt7\)
\end{answer}

\section{解答题}

\begin{problem}
在 \(\triangle ABC\) 中，内角 \(A\)，\(B\)，\(C\) 所对的边分别为 \(a\)，\(b\)，\(c\)，已知 \(b+c=2a\cos B\).

\begin{enumerate}
\item 证明：\(A=2B\)；
\item 若 \(\cos B=\dfrac23\)，求 \(\cos C\) 的值.
\end{enumerate}
\end{problem}

\begin{solution}
\begin{enumerate}
\item 由正弦定理，设三角形 \(ABC\) 的外接圆半径为 \(R\)，则
\[
b+c=2R(\sin B+\sin C),\qquad a=2R\sin A.
\]
因此由已知条件得 \(\sin B+\sin C=2\sin A\cos B\). 又因为 \(b+c>0\)，所以 \(\cos B>0\)，从而 \(B\in(0,\frac{\pi}{2})\). 由 \(C=\pi-A-B\)，有 \(\sin C=\sin(A+B)\)，故
\[
\sin A\cos B+(1+\cos A)\sin B=2\sin A\cos B,
\]
即 \((1+\cos A)\sin B=\sin A\cos B\). 由于 \(A\in(0,\pi)\)，且 \(B\in(0,\frac{\pi}{2})\)，可得
\[
\tan B=\frac{\sin A}{1+\cos A}=\tan\frac A2.
\]
又 \(B\) 与 \(\frac A2\) 均在 \((0,\frac{\pi}{2})\) 内，正切函数在此区间上单调递增，所以 \(B=\frac A2\)，即 \(A=2B\).

\item 由第（1）问，\(C=\pi-A-B=\pi-3B\). 因为 \(\cos B=\frac23\)，所以
\[
\cos C=-\cos 3B=-\left(4\cos^3B-3\cos B\right)
=-\left(4\left(\frac23\right)^3-3\cdot\frac23\right)=\frac{22}{27}.
\]
\end{enumerate}
\end{solution}

\begin{problem}
设数列 \(\{a_n\}\) 的前 \(n\) 项和为 \(S_n\)，已知 \(S_2=4\)，\(a_{n+1}=2S_n+1\)，\(n\in\N^*\).

\begin{enumerate}
\item 求通项公式 \(a_n\)；
\item 求数列 \(\{|a_n-n-2|\}\) 的前 \(n\) 项和.
\end{enumerate}
\end{problem}

\begin{solution}
\begin{enumerate}
\item 记 \(S_1=a_1\). 由 \(a_2=2S_1+1\) 及 \(S_2=S_1+a_2=4\)，得
\[
3S_1+1=4,\qquad S_1=a_1=1.
\]
对任意 \(n\in\N^*\)，由递推关系得
\[
S_{n+1}=S_n+a_{n+1}=3S_n+1.
\]
于是 \(S_{n+1}+\frac12=3\left(S_n+\frac12\right)\). 结合 \(S_1=1\)，可得
\[
S_n+\frac12=3^{n-1}\left(S_1+\frac12\right)=\frac{3^n}{2},
\]
即 \(S_n=\frac{3^n-1}{2}\). 因此 \(a_1=1\)，且对 \(n\geqslant2\)，
\[
a_n=S_n-S_{n-1}=\frac{3^n-1}{2}-\frac{3^{n-1}-1}{2}=3^{n-1}.
\]
故对一切 \(n\in\N^*\)，\(a_n=3^{n-1}\).

\item 由 \(a_1=1\)，\(a_2=3\)，可得
\[
|a_1-1-2|=2,\qquad |a_2-2-2|=1.
\]
当 \(k\geqslant3\) 时，\(3^{k-1}>k+2\)：当 \(k=3\) 时不等式成立；若对某个 \(k\geqslant3\) 成立，则 \(3^k>3(k+2)>k+3\)，所以可由归纳法得证. 因此
\[
|a_k-k-2|=3^{k-1}-k-2\quad(k\geqslant3).
\]
当 \(n=1\) 或 \(n=2\) 时，前 \(n\) 项和分别为 \(2\) 和 \(3\). 当 \(n\geqslant3\) 时，
\[
\sum_{k=1}^{n}|a_k-k-2|=3+\frac{3^n-9}{2}-\left(\frac{n(n+1)}2+2n-7\right).
\]
化简得
\[
\sum_{k=1}^{n}|a_k-k-2|=\frac{3^n-n^2-5n+11}{2}.
\]
综上，数列 \(\{|a_n-n-2|\}\) 的前 \(n\) 项和为
\[
\sum_{k=1}^{n}|a_k-k-2|=
\begin{cases}
2, & n=1,\\
3, & n=2,\\
\dfrac{3^n-n^2-5n+11}{2}, & n\geqslant3.
\end{cases}
\]
\end{enumerate}
\end{solution}

\begin{problem}
如图，在三棱台 \(ABC-DEF\) 中，平面 \(BCFE\perp\) 平面 \(ABC\)，\(\angle ACB=90^\circ\)，\(BE=EF=FC=1\)，\(BC=2\)，\(AC=3\).

\begin{enumerate}
\item 求证：\(BF\perp\) 平面 \(ACFD\)；
\item 求直线 \(BD\) 与平面 \(ACFD\) 所成角的余弦值.
\end{enumerate}

\bitmapfigure[float=inline,max width=9.3cm]{img/2016/zhejiang_liberal/q18_fig1.png}
\end{problem}

\begin{solution}
建立如图所示的直角坐标系，使平面 \(ABC\) 为 \(xOy\) 平面，平面 \(BCFE\) 为 \(xOz\) 平面，并取
\[
C=(0,0,0),\qquad B=(2,0,0),\qquad A=(0,3,0).
\]
因为三棱台的对应底边平行，所以 \(EF\parallel BC\)，且 \(\dfrac{EF}{BC}=\frac12\). 设 \(E=(\frac32,0,h)\)，\(F=(\frac12,0,h)\). 由 \(BE=1\) 得
\[
h^2+\left(2-\frac32\right)^2=1,\qquad h=\frac{\sqrt3}{2}.
\]
又因为 \(DF\parallel AC\) 且 \(DF=\frac12AC\)，所以
\[
D=\left(\frac12,\frac32,\frac{\sqrt3}{2}\right).
\]

\begin{enumerate}
\item 平面 \(ACFD\) 内的两条相交直线可取 \(CA\) 和 \(CF\). 有
\[
\overrightarrow{BF}=\left(-\frac32,0,\frac{\sqrt3}{2}\right),\quad
\overrightarrow{CA}=(0,3,0),\quad
\overrightarrow{CF}=\left(\frac12,0,\frac{\sqrt3}{2}\right).
\]
于是
\[
\overrightarrow{BF}\cdot\overrightarrow{CA}=0,\qquad
\overrightarrow{BF}\cdot\overrightarrow{CF}=-\frac34+\frac34=0.
\]
所以 \(BF\) 同时垂直于平面 \(ACFD\) 内相交的两条直线 \(CA\) 和 \(CF\)，故 \(BF\perp\) 平面 \(ACFD\).

\item 设直线 \(BD\) 与平面 \(ACFD\) 所成的角为 \(\theta\). 由第（1）问，\(BF\) 是平面 \(ACFD\) 的法线，因此
\[
\sin\theta=\frac{|\overrightarrow{BD}\cdot\overrightarrow{BF}|}{|\overrightarrow{BD}|\,|\overrightarrow{BF}|}.
\]
其中
\[
\overrightarrow{BD}=\left(-\frac32,\frac32,\frac{\sqrt3}{2}\right),
\quad |\overrightarrow{BD}|=\frac{\sqrt{21}}2,
\quad |\overrightarrow{BF}|=\sqrt3,
\]
且
\[
\overrightarrow{BD}\cdot\overrightarrow{BF}=\frac94+\frac34=3.
\]
所以 \(\sin\theta=\frac{2}{\sqrt7}\)，从而
\[
\cos\theta=\sqrt{1-\sin^2\theta}=\sqrt{1-\frac47}=\frac{\sqrt{21}}7.
\]
\end{enumerate}
\end{solution}

\begin{problem}
如图，设抛物线 \(y^2=2px\)（\(p>0\)）的焦点为 \(F\)，抛物线上的点 \(A\) 到 \(y\) 轴的距离等于 \(|AF|-1\).

\begin{enumerate}
\item 求 \(p\) 的值；
\item 若直线 \(AF\) 交抛物线于另一点 \(B\)，过 \(B\) 与 \(x\) 轴平行的直线和过 \(F\) 与 \(AB\) 垂直的直线交于点 \(N\)，\(AN\) 与 \(x\) 轴交于点 \(M\)，求 \(M\) 的横坐标的取值范围.
\end{enumerate}

\FigureLayoutDeclare{img/2016/zhejiang_liberal/q19_fig1.png}{6.0cm}{85bp 110bp 222bp 83bp}
\bitmapinlinefigure[7.0cm][max width=7.0cm]{img/2016/zhejiang_liberal/q19_fig1.png}
\end{problem}

\begin{solution}
\begin{enumerate}
\item 设抛物线上的点 \(A\) 的参数为 \(t\)，则
\[
A=\left(\frac p2t^2,pt\right),\qquad F=\left(\frac p2,0\right).
\]
由于抛物线在右侧，点 \(A\) 到 \(y\) 轴的距离为 \(\frac p2t^2\)，而
\[
|AF|=\sqrt{\left(\frac p2t^2-\frac p2\right)^2+p^2t^2}=\frac p2(t^2+1).
\]
由题意，\(\frac p2t^2=\frac p2(t^2+1)-1\)，所以 \(p=2\).

\item 此时抛物线方程为 \(y^2=4x\)，焦点为 \(F=(1,0)\)，且
\[
A=(t^2,2t).
\]
设直线 \(AF\) 与抛物线的另一交点为 \(B\)，记 \(B\) 的参数为 \(s\). 过参数为 \(t,s\) 的两点的弦所在直线方程为
\[
(t+s)y=2x+2ts.
\]
因为焦点 \(F=(1,0)\) 在直线 \(AB\) 上，所以 \(2+2ts=0\)，即 \(ts=-1\)，从而
\[
B=\left(\frac1{t^2},-\frac2t\right),\qquad t\ne0.
\]
过 \(B\) 的水平直线为 \(y=-\frac2t\). 直线 \(AB\) 的方向向量可取 \((t^2-1,2t)\). 设过 \(F\) 且垂直于 \(AB\) 的直线与上述水平直线交于 \(N=(u,-\frac2t)\)，则
\[
(u-1,-\tfrac2t)\cdot(t^2-1,2t)=0,
\]
所以 \(u=1+\frac4{t^2-1}\).

当 \(t^2=1\) 时，\(AB\) 为过 \(F\) 的竖直直线，过 \(F\) 的垂线与过 \(B\) 的水平线平行，点 \(N\) 不存在，故还需有 \(t^2\ne1\).

设 \(M=(m,0)\). 因为 \(A,N,M\) 三点共线，由截距关系得
\[
m=\frac{t^2(1+u)}{t^2+1}=\frac{2t^2}{t^2-1}.
\]
令 \(v=t^2\)，则 \(v>0\) 且 \(v\ne1\). 当 \(0<v<1\) 时，\(m=\frac{2v}{v-1}\) 的取值范围为 \((-\infty,0)\)；当 \(v>1\) 时，\(m=2+\frac2{v-1}\) 的取值范围为 \((2,+\infty)\). 因此 \(M\) 的横坐标的取值范围为
\[
(-\infty,0)\cup(2,+\infty).
\]
\end{enumerate}
\end{solution}

\begin{problem}
设函数 \(f(x)=x^3+\dfrac{1}{x+1}\)，\(x\in[0,1]\)，证明：

\begin{enumerate}
\item \(f(x)\ge1-x+x^2\)；
\item \(\dfrac34<f(x)\le\dfrac32\).
\end{enumerate}
\end{problem}

\begin{solution}
\begin{enumerate}
\item 在 \(x\in[0,1]\) 时，\(x+1>0\)，于是
\[
f(x)-(1-x+x^2)=x^3+\frac1{x+1}-1+x-x^2=\frac{x^4}{x+1}\geqslant0.
\]
所以 \(f(x)\geqslant1-x+x^2\).

\item 由
\[
1-x+x^2=\left(x-\frac12\right)^2+\frac34\geqslant\frac34
\]
及第（1）问知 \(f(x)\geqslant\frac34\). 若 \(x\ne\frac12\)，则上式中的二次式严格大于 \(\frac34\)；若 \(x=\frac12\)，则
\[
f(x)-(1-x+x^2)=\frac{x^4}{x+1}>0.
\]
故始终有 \(f(x)>\frac34\).

另一方面，
\[
\frac32-f(x)=\frac32-x^3-\frac1{x+1}=\frac{(1-x)(2x^3+4x^2+4x+1)}{2(x+1)}\geqslant0,
\]
其中最后一个不等号由 \(x\in[0,1]\) 得到. 所以 \(f(x)\leqslant\frac32\)，且当 \(x=1\) 时取等号. 综上，
\[
\frac34<f(x)\leqslant\frac32.
\]
\end{enumerate}
\end{solution}
